% Capitulo: Arquitetura medalhão e modelagem tradicional
\chapter{Arquitetura medalhão e modelagem tradicional}
\label{cap:medalhao}

\rascunho{
  \textbf{Ideia do capítulo.} Abrir a Parte~\ref{parte:modelagem} explicando a
  arquitetura medalhão (bronze, prata, ouro), por que ela ficou popular na era do
  lakehouse e em que ela é diferente da modelagem tradicional para análise descritiva.
  A tese: medalhão organiza camadas de refinamento, não substitui modelagem. E o mesmo
  dado, bem organizado em camadas, passa a servir a usos que a modelagem tradicional
  não previa: ciência de dados, treino de IA e contexto para IA.
  \begin{itemize}
    \item Retomar Inmon e Kimball do capítulo~\ref{cap:dw} e mostrar onde cada um
      reaparece dentro do medalhão.
    \item A Figura~\ref{fig:medalhao} e a Tabela~\ref{tab:usos-modelagem} servem de mapa
      para os próximos três capítulos.
  \end{itemize}
}

% ---------------------------------------------------------------------------
\section{Bronze, prata e ouro}
\label{sec:bronze-prata-ouro}

\rascunho{
  \begin{itemize}
    \item O termo vem do ecossistema do lakehouse \parencite{databricksmedallion,armbrust2021}.
    \item Bronze: o dado como chegou da fonte, sem transformação, com histórico. Serve
      para reprocessar e auditar.
    \item Prata: limpo, tipado, sem duplicidade, com as chaves integradas entre fontes.
    \item Ouro: pronto para um consumo específico (métricas de negócio, modelos
      dimensionais, tabelas de features).
    \item O paralelo com as camadas de um projeto de analytics engineering (preparação,
      intermediária, marts). \pendente{decidir se uso a nomenclatura de alguma ferramenta como exemplo}
  \end{itemize}
}

\pendente{escrever}

\begin{figure}[htbp]
  \centering
  % Arquitetura medalhao e os consumidores de cada camada (cap. 12).
\begin{tikzpicture}[diagrama, x=1mm, y=1mm]
  % fontes
  \foreach \y/\nome/\id in {12/Sistemas/f1, 4/APIs/f2, -4/Arquivos/f3, -12/Documentos/f4}
    \node[caixa, minimum width=15mm, minimum height=6mm] (\id) at (7.5,\y) {\nome};
  \node[moldura, fit=(f1)(f4), label={[rotulo]above:Fontes}] {};

  % camadas
  \node[banco, minimum width=19mm, minimum height=26mm] (b) at (31,0)
    {\textbf{Bronze}\\[1mm]{\scriptsize bruto, como}\\{\scriptsize chegou}};
  \node[banco, minimum width=19mm, minimum height=26mm] (p) at (53,0)
    {\textbf{Prata}\\[1mm]{\scriptsize limpo, tipado,}\\{\scriptsize integrado}};
  \node[bancodestaque, minimum width=19mm, minimum height=26mm] (o) at (75,0)
    {\textbf{Ouro}\\[1mm]{\scriptsize pronto para}\\{\scriptsize o consumo}};
  \foreach \id in {f1,f2,f3,f4} \draw[seta] (\id.east) -- (b.west);
  \draw[seta] (b) -- (p);
  \draw[seta] (p) -- (o);

  % consumidores
  \foreach \y/\nome/\id in {15/{BI descritivo}/c1, 5/{Ciência de dados}/c2, -5/{Treino de IA}/c3, -15/{Contexto p/ IA}/c4}
    \node[caixa, minimum width=21mm, minimum height=6mm] (\id) at (104.5,\y) {\nome};
  \foreach \id in {c1,c2,c3,c4} \draw[seta] (o.east) -- (\id.west);
  \node[rotulo, text width=110mm] at (58,-23)
    {Ciência de dados e treino de IA muitas vezes leem também da prata, que guarda mais detalhe que o ouro.};
\end{tikzpicture}

  \caption{A arquitetura medalhão e os consumidores que leem de cada camada.}
  \label{fig:medalhao}
  \fonte{elaborado pelo autor com base em \textcite{databricksmedallion}.}
\end{figure}

% ---------------------------------------------------------------------------
\section{O que muda em relação a Inmon e Kimball}
\label{sec:medalhao-tradicional}

\rascunho{
  \begin{itemize}
    \item Medalhão fala de \emph{camadas de refinamento}; Inmon e Kimball falam de
      \emph{modelo}. Por isso não são concorrentes diretos.
    \item A prata lembra a camada integrada de Inmon; o ouro, quando serve a BI, quase
      sempre é um modelo dimensional de Kimball (ver \veja{sec:modelagem-dimensional}).
    \item O que mudou foi o custo: com armazenamento barato, guardar o bruto (bronze) e
      várias versões refinadas deixou de ser luxo.
    \item O risco: tratar medalhão como receita e acabar com três cópias do mesmo dado e
      nenhum modelo. Bronze, prata e ouro sem grão definido continuam sendo bagunça, só
      que organizada em pastas.
    \item Relato: como eu organizei camadas num projeto. \pendente{escolher o caso}
  \end{itemize}
}

\pendente{escrever}

% ---------------------------------------------------------------------------
\section{Modelagem para análise descritiva}
\label{sec:modelagem-descritiva}

\rascunho{
  \begin{itemize}
    \item A análise descritiva responde o que aconteceu e por quê. É o uso que a
      modelagem tradicional sempre atendeu, e continua sendo o maior volume de trabalho
      na maioria das empresas.
    \item Na camada ouro: esquema estrela, métricas definidas uma vez, tabelas agregadas
      para os painéis mais acessados.
    \item O que vale manter da tradição: grão explícito, dimensões conformadas, chaves
      substitutas, histórico das dimensões.
  \end{itemize}
}

\pendente{escrever}

% ---------------------------------------------------------------------------
\section{Um dado, vários usos}
\label{sec:um-dado-varios-usos}

\rascunho{
  \begin{itemize}
    \item A modelagem tradicional tinha um consumidor principal: o analista com a
      ferramenta de BI. Hoje o mesmo dado alimenta modelos preditivos, decisões
      automatizadas, treino de IA e assistentes que respondem perguntas.
    \item Cada uso pede uma forma diferente de dado (Tabela~\ref{tab:usos-modelagem}).
      Os próximos capítulos tratam um uso cada.
  \end{itemize}
}

\pendente{escrever}

\begin{table}[htbp]
  \centering
  \caption{Cada uso pede uma forma diferente de organizar os dados.}
  \label{tab:usos-modelagem}
  \small
  \begin{tabularx}{\textwidth}{@{}>{\raggedright\arraybackslash}p{0.22\textwidth}
      >{\raggedright\arraybackslash}X>{\raggedright\arraybackslash}X@{}}
    \toprule
    Uso & Forma dos dados & Cuidado principal \\
    \midrule
    Descritiva \newline{\footnotesize\itshape o que aconteceu?} & Esquema estrela e tabelas agregadas, na camada ouro & Métrica definida uma vez; grão explícito \\
    \addlinespace
    Preditiva \newline{\footnotesize\itshape o que vai acontecer?} & Tabela larga: uma linha por entidade e momento, com variáveis e alvo & Não deixar o futuro vazar para o passado \\
    \addlinespace
    Prescritiva \newline{\footnotesize\itshape o que fazer?} & Previsões junto com custos, restrições e cenários do negócio & Registrar a decisão tomada e o resultado \\
    \addlinespace
    Treino de IA & Conjunto curado, rotulado, versionado e documentado & Qualidade e representatividade dos exemplos, privacidade \\
    \addlinespace
    Contexto para IA & Documentos em trechos, metadados, glossário, camada semântica & Atualização, dono de cada fonte, permissões \\
    \bottomrule
  \end{tabularx}
  \fonte{elaborado pelo autor.}
\end{table}

% ---------------------------------------------------------------------------
\begin{leituras}
  \leitura{armbrust2021}{O artigo que propõe o lakehouse, o contexto em que a
    arquitetura medalhão se popularizou.}
  \leitura{kleppmann2017}{Para entender os sistemas por baixo das camadas:
    armazenamento, replicação, lote e fluxo contínuo, sem depender de ferramenta.}
\end{leituras}
